\begin{proposition}[C.3.5 (Known formula, with proved comparison): the closing translation]
Fix the standard twists and all root and cusp markings of R2.
There is a coherent choice of the even-shift maps
\(\widetilde F_{h,c}:X_c\to X_{c+h}\), \(h\in2\mathbb Z\),
with the following explicit ordinary-chart translation. Choose global
roots \(F^3=t\), \(G^4=t-1\) on the uniformizing cover and put
\(\ell=(2\pi i)^{-1}\log(F/G)\pmod{\mathbb Z}\), \(e=(0,1)^t\).
Then the translation is
\begin{equation}\label{eq:C-explicit-closing}
 b_{h,c}=-\frac h{12}\Pi_c\widehat\gamma
              +\left(h\ell-\frac{h^2}{24}\right)e.
\end{equation}
For \(h=12\), it is the scalar multiplication
\begin{equation}\label{eq:C-twelve-scalar}
 K(t)=\frac{t^4}{(t-1)^3}
\end{equation}
on the ordinary \(\mathbb C^*\)-coordinate, after removing an integral
period. On the finite root charts its scalar representatives are
\(K(t)s_1^{-12}\) and \(K(t)s_2^{12}\), respectively, both units.
In a cusp coordinate \(q=t_c\), the scalar representative is
\(q^{-1}H(q)\), with \(H\) a holomorphic nonvanishing function.
The same map has a holomorphic \((\mathbb C^*)^2\)-unit representative
on the entire cusp filling. These two representatives differ by
exactly the \(\widehat\gamma\) cusp deck transformation.

For six positive \(w\)-loops the closing map is
\(\widetilde F_{12,c_0-12}\); thus the stated cusp exponent is \(-1\)
with the orientation convention of C.1.1. This computes the closing
cocycle, not the mapping class of its composition with parallel
transport. In particular it does not imply \(6[\Phi]=0\).
\end{proposition}
\begin{proof}
The raw formula without \(-h^2e/24\), its three extension checks,
and the coherent correction occur in
\cite[\texttt{p2:prop:global-even-shift} and the subsequent
orbit-groupoid calculations]{DGKXv1push}. We verify their exact use
here; the source is version~1, and version~2 omits these statements.
The global roots exist because the zero orders of \(t\) and \(t-1\)
on the simply connected uniformizing surface are divisible by three
and four. The local characters and their coprime orders give
\[
 \ell(g_1z)-\ell(z)=-\frac13,\qquad
 \ell(g_2z)-\ell(z)=\frac14,\qquad
 \ell(g_0z)-\ell(z)=\frac1{12}\pmod{\mathbb Z}.
\]
A change of logarithm adds an integral \(h\widehat\delta\) period,
so \eqref{eq:C-explicit-closing} is a well-defined torus translation.
Since \(\Pi_{c+h}=\Pi_c(I+hE)\), the two ordinary lattices agree;
the source-to-target lattice map of the identity linear part is
\(I-hE\).

For compatible logarithm lifts, covariance gives the affine defects
\[
 b_{h,c}(g_jz)-R_j(z)b_{h,c}(z)
       =\Pi_c(g_jz)\eta_j,\qquad
 \eta_1=\frac h2(\widehat u-\widehat w-\widehat\delta),\quad
 \eta_2=\frac h2(-\widehat w+\widehat\delta).
\]
Their \(\gamma\)-coordinates vanish, so the same right sides use
\(\Pi_{c+h}\). The vectors are integral for even \(h\). At the
cusp the defect is zero: \((A_0-I)\widehat\gamma=-\widehat\delta\)
cancels the increment of \(\ell\). These equations prove ordinary
and cusp descent with the stated orientation of the peripheral
loops.

For the finite comparison put \(w_j=(2\pi i)^{-1}\log s_j\).
The functions \(F/(s_1G)\) and \(Fs_2/G\) are invariant
holomorphic units on the respective full root discs. Choose
holomorphic logarithms \(\kappa_j\) of them. Then
\[
 \ell=\epsilon_jw_j+\kappa_j\pmod{\mathbb Z},\qquad
 \epsilon_1=1,\quad\epsilon_2=-1.
\]
The exact finite-to-ordinary comparison satisfies
\(\sigma_{j,c+h}-\sigma_{j,c}=h\epsilon_jw_je\).
Consequently the required finite-chart translation is
\[
 b^{(j)}_{h,c}
 =-\frac h{12}\Pi_c\widehat\gamma
       +\left(h\kappa_j-\frac{h^2}{24}\right)e.
\]
It is holomorphic on the full root disc and satisfies
\[
 b^{(j)}_{h,c}(g_jz)-R_j(z)b^{(j)}_{h,c}(z)
       -\frac{h\epsilon_j}{m_j}e
 =\Pi_c(g_jz)\eta_j.
\]
This is precisely equivariance for the source and target affine
finite actions, including the changed period column. The exact
identity
\(\sigma_{j,c+h}+b^{(j)}_{h,c}
=\sigma_{j,c}+b_{h,c}\) proves overlap agreement. Thus no real
interpolation of a multivalued period has been used.

At the cusp let \(q=e^{2\pi i s}\) and
\(C_{c,21}=\beta^{\rm part}+c+s\), which is holomorphic at
\(q=0\). The function
\[
 \widehat H=e^{2\pi i s/12}F/G
 \quad\hbox{satisfies}\quad
 \widehat H^{12}=q\frac{t^4}{(t-1)^3}.
\]
The right side is a unit because \(t\) has a simple pole at the
cusp. The character formula makes \(\widehat H\) single-valued;
it therefore extends as a nonvanishing holomorphic function. Write
\(\ell=-s/12+\kappa_0(q)\), with \(\kappa_0\) holomorphic.
Since
\(\Pi_c\widehat\gamma=(6\mu,-s+C_{c,21})^t\),
the logarithmic terms cancel in \eqref{eq:C-explicit-closing}:
\[
 \bar b_{h,c}(q)=
 \left(-\frac h2\mu(q),\,
 -\frac h{12}C_{c,21}(q)+h\kappa_0(q)-\frac{h^2}{24}\right)^t.
\]
Exponentiating gives a holomorphic unit in \((\mathbb C^*)^2\).
On every toric chart its action multiplies a character coordinate
by an integral character of this unit, so it and its inverse
extend across every double curve and both triple points. The
source and target cusp lattice multipliers differ by
\((1,e^{2\pi i h\lambda_1})=(1,1)\), for
\(\lambda\in\mathbb Z^2\). They are therefore identical. The
unit multiplication commutes with their toric shears and descends.
Exponentiation of the displayed formula gives exactly the ordinary
translation on the punctured cusp. We have now verified all
three comparisons and their inverses.

The uncorrected translations have composition defect
\(-hk e/12\). The constants in
\eqref{eq:C-explicit-closing} cancel it, since
\((h+k)^2-h^2-k^2=2hk\). Hence
\[
 \widetilde F_{k,c+h}\widetilde F_{h,c}
       =\widetilde F_{h+k,c},\qquad
 \widetilde F_{0,c}=\mathrm{id}.
\]
Any difference from the fixed R2 choice is a scalar automorphism,
which lies in its connected \(\mathbb C^*\). The choice therefore
has the same smooth monodromy class.

For \(h=12\), remove the actual period
\(-\Pi_c\widehat\gamma\) from \eqref{eq:C-explicit-closing}.
The remaining translation is \((12\ell-6)e\). Its exponential
is \((F/G)^{12}=K(t)\), proving
\eqref{eq:C-twelve-scalar}. The finite corrections divide it by
\(s_1^{12}\) and multiply it by \(s_2^{12}\). The original
uniformizers are retained: these are units because
\(t/s_1^3\) and \((t-1)/s_2^4\) are units. They need not be
replaced by one.

Put \(H(q)=qK(t(q))=e^{24\pi i\kappa_0(q)}\). The cusp unit
for \(h=12\) is exactly
\[
 U_{12,c}(q)=
 \left(e^{-12\pi i\mu(q)},\,
       e^{-2\pi i C_{c,21}(q)}H(q)\right).
\]
The \(\widehat\gamma\) deck transformation has multiplier
\[
 \left(e^{12\pi i\mu(q)},\,
       e^{2\pi i C_{c,21}(q)}q^{-1}\right),
\]
because \(B_0(1,0)^t=(0,-1)^t\). Multiplying these two
expressions gives \((1,q^{-1}H)=(1,K)\). This proves the
claimed comparison, including its sign, on the entire cusp
quotient. In particular the scalar representative's pole is not
an obstruction to the holomorphic map itself: its unit
representative already extends.

Finally \(w=e^{-\pi i c}\) makes six positive loops the path
\(c(r)=c_0-12r\); its endpoint is closed by the positive shift
\(h=12\). A quotient isotopy proving an order bound would still
have to patch the ordinary scalar description with the two-coordinate
cusp unit description. Scaling their additive logarithms separately
produces on their overlap a fractional
\(\widehat\gamma\)-period, which is not an integral period during
the interpolation. This patching has not been established. The
exact substitutions, coherent correction, finite defects, and
cusp comparison are checked in
\path{code/d_push_monodromy.py} and
\path{code/d_push_closing.py}; these are symbolic checks of the
identities proved above.
\end{proof}
